\section{Monistot}
Viite \citep{Neumann1989}

Käyrien muodostamien linkkien tutkimiseen tarvitaan (sileiden) monistojen teoriaa. Monistot ovat geometrisia rakenteita, jotka muistuttavat paikallisesti euklidista avaruutta. Ne mahdollistavat \emph{avaruuden rakenteen} käsitteen. Algebrallisten solmujen kannalta tämä on olennaista, sillä solmujen tutkimista helpottaa niiden upottaminen monistoon. Tässä työssä noudatetaan lähteen \citep{Lee2012} merkintätapoja.
\begin{maaritelma}
    Olkoon $M$ topologinen avaruus. $M$ on \emph{$n$-ulotteinen topologinen monisto}, jos
    \begin{itemize}
        \item $M$ on Hausdorff-avaruus
        \item $M$ on $N_2$-avaruus
        \item $M$ on paikallisesti euklidisesti $n$-ulotteinen: kaikilla $p\in M$ on avoin ympäristö $U\subset M$, joka on homeomorfinen johonkin $\R^n$:n avoimeen osajoukkoon $V\subset\R^n$.
    \end{itemize}
\end{maaritelma}
\begin{huomautus}
    Yleisesti kirjoitetaan ``olkoon $M^n$ monisto'', jolla tarkoitetaan, että $M$ on $n$-ulotteinen monisto.
\end{huomautus}
\begin{esimerkki}
    Esimerkiksi Maapallo on 2-ulotteinen monisto, koska paikallisesti Maapallon pinta muistuttaa tasoa. Maapallo on erityistapaus algebrallisesta pinnasta, koska sen polynomiyhtälö $x^2+y^2+z^2=R^2$.
\end{esimerkki}
\begin{maaritelma}
    Olkoon $M^n$ monisto. $M$:n \emph{koordinaattikartta} on pari $(U,\varphi)$, jossa $U\subset M$ on avoin ja $\varphi:U\to\hat U$ on homeomorfismi avoimeen joukkoon $\hat U\subset R^n$. Koordinaattikartan \emph{koordinaattikuvaus} on $\varphi$.
\end{maaritelma}
\begin{esimerkki}
    Yleisesti algebrallinen hyperpinta $f(x_1,x_2,\dots,x_{n+1})=0$ määrittää $n$-ulotteisen moniston $\R^{n+1}$:ssä.
\end{esimerkki}
\begin{maaritelma}
    Jos $U\subset\R^n,V\subset\R^m$, funktio $f:U\to V$ on \emph{sileä}, jos sen jokaisella komponenttifunktiolla on jatkuvat asteen $n$ osittaisderivaatat kaikilla $n\in\N$.
\end{maaritelma}
\begin{maaritelma}
    Kuvaus $f:A\to B$ on \emph{diffeomorfismi}, jos $f$ on bijektio ja $f$ ja $f^{-1}$ ovat sileitä.
\end{maaritelma}
\begin{maaritelma}
    Moniston $M$ \emph{atlas} $\mathcal{A}$ on kokoelma karttoja $\{(U_i,\varphi_i)\}_{i\in I}$ siten, että karttojen lähtöjoukot peittävät moniston $M$: kaikille $p\in M$ on olemassa $i\in I$, jolle $p\in U_i$. Jos koordinaattikartat ovat pareittain sileästi yhteensopivia, eli
    \begin{equation}
        \varphi_j\circ\varphi_i^{-1}:\varphi_i(U_i\cap U_j)\to\varphi_j(U_i\cap U_j)
    \end{equation}
    tai $U_i\cap U_j=\emptyset$, niin $\mathcal{A}$ on \emph{sileä atlas}.
\end{maaritelma}
\begin{maaritelma}
    Monisto $M$ on \emph{sileä monisto}, jos sillä on sileä atlas.
\end{maaritelma}


Tässä päätulos:
\begin{lause}
    Pelkistetty algebrallinen käyrä $V\subset\C^2$ valitulla upotuksella $\C^2\subset\mathbb{P}^2$ määrittää RPI-pujoskaavion $\Omega$ käyrän $V$ sidokselle äärettömyydessä. Pisteiden lukumäärä äärettömyydessä $n_\infty(V)$ on $\Omega$:n juurisolmun valenssi $n$ ja $deg(V)$ on juurisolmun ja sen sidoksen virtuaalisen komponentin sidosluku $s_\text{juuri}$.
\end{lause}
\begin{lause}
    Sidos on jonkin äärettömyydessä olevan sidoksen (joka voidaan valita säännölliseksi tai jopa hyväksi) osasidos jos ja vain jos sillä on RPI-pujoskaavio.
\end{lause}


\subsection{Topologiset ja algebralliset määritelmät}
\begin{maaritelma}
    Linkki äärettömyydessä
\end{maaritelma}
\begin{maaritelma}
    Kuitu, säännöllinen kuitu
\end{maaritelma}
\begin{maaritelma}
    Säännöllisyys äärettömyydessä
\end{maaritelma}
\begin{maaritelma}
    Milnorin kidutus
\end{maaritelma}
\begin{maaritelma}
    Säännöllinen algebrallinen käyrä
\end{maaritelma}
\begin{propositio}
    Invariantit
\end{propositio}
\subsection{Kaapelien ja kaavioiden teoria}
\begin{maaritelma}
    Torussidos
\end{maaritelma}
\begin{maaritelma}
    Pujoskaavio
\end{maaritelma}
\begin{maaritelma}
    Kaapelointi tai säikeistys
\end{maaritelma}
\begin{maaritelma}
    Juurisolmu
\end{maaritelma}
\begin{lause}
    Käänteiset Puiseux-epäyhtälöt
\end{lause}